Accurate prediction of spatio-temporal responses in geotechnical systems remains challenging due to the high computational cost of numerical simulations and the coupled multi-physics nature of soil behaviour. This paper presents an engineering-friendly neural network surrogate modelling framework for predicting time-dependent, multi-output responses from static geotechnical input parameters. The proposed approach provides a practical and unified surrogate representation in which the temporal evolution of the system is captured within a shared latent structure, while different spatial and multi-physics responses are generated through output-specific reconstruction modules using a simplified architecture that focuses on temporal modelling and spatial reconstruction, resulting in straightforward deployment in engineering practice. This formulation enables multiple responses, including scalar indicators, vector quantities and high-dimensional spatial fields such as stress and pore pressure distributions, to be predicted consistently within a single surrogate model, avoiding the need to construct and manage separate models for individual outputs or response fields. The framework is computationally efficient and suitable for the repeated evaluations required in applications such as parametric studies and uncertainty analysis. The approach is assessed using two benchmark problems of increasing complexity, including an offshore strip footing loading case and a sequential triaxial test with pronounced nonlinearity and a higher-dimensional input space. The proposed metamodel provides an effective and practically accurate representation of the spatiotemporal response: the test-set global normalised absolute error (GNAE) is about 2.3\% in the offshore footing example and about 1.3\% in the triaxial test. These results demonstrate that the surrogate can reproduce complex soil responses with good accuracy, highlighting its potential as a practical tool for spatio-temporal surrogate modelling in geotechnical engineering practice.



\keywords{Surrogate modelling; Neural network; Decoder; Spatio-temporal; Multi-field; Geotechnical engineering; Sequential inversion}



